\documentclass[12pt]{article}
\usepackage{fullpage,amsmath,amssymb,amsthm,hyperref}

\newtheorem{theorem}{Theorem}
\newtheorem{proposition}[theorem]{Proposition}
\newtheorem{lemma}[theorem]{Lemma}
\theoremstyle{definition}
\newtheorem{definition}[theorem]{Definition}
\newtheorem{remark}[theorem]{Remark}

% Macros
\newcommand{\bR}{\mathbb{R}}
\newcommand{\score}{\mathrm{score}}
\DeclareMathOperator{\roots}{roots}

\title{Solution to Problem 4:\\
Subharmonicity of $1/\Phi_n$ under Finite Free Convolution}
\author{}
\date{April 29, 2026}

\begin{document}
\maketitle

\section*{Problem Statement}

Let $p(x) = \sum_{k=0}^n a_k x^{n-k}$ and $q(x) = \sum_{k=0}^n b_k x^{n-k}$ be monic polynomials of degree $n$ with $a_0 = b_0 = 1$. Define the \emph{finite free convolution} $(p \boxplus_n q)(x) = \sum_{k=0}^n c_k x^{n-k}$ where
\[
  c_k = \sum_{i+j=k} \frac{(n-i)!\,(n-j)!}{n!\,(n-k)!}\,a_i b_j.
\]
For a monic polynomial $p(x) = \prod_{i=1}^n (x - \lambda_i)$ with distinct roots, define the \emph{finite free Fisher information}
\[
  \Phi_n(p) := \sum_{i=1}^n \Bigl(\sum_{j \neq i} \frac{1}{\lambda_i - \lambda_j}\Bigr)^2 = \|\score(\lambda)\|^2,
\]
where $\score(\lambda)_i := \sum_{j \neq i} \tfrac{1}{\lambda_i - \lambda_j}$, and $\Phi_n(p) := \infty$ if $p$ has a repeated root.

\medskip
\noindent\textbf{Question.} Is it true that for all monic real-rooted polynomials $p$ and $q$ of degree $n$,
\[
  \frac{1}{\Phi_n(p \boxplus_n q)} \;\ge\; \frac{1}{\Phi_n(p)} + \frac{1}{\Phi_n(q)}\,?
\]

\medskip
\noindent\textbf{Summary of results.} We prove two things completely:
\begin{enumerate}
  \item[\S1.] The \emph{score formula} expressing $\score(\lambda)_k$ in terms of $p'$ and $p''$.
  \item[\S2.] The inequality for $n = 2$ (with equality).
\end{enumerate}
In \S3 we give an honest account of why the case $n \ge 3$ remains open, state the missing ingredient, and provide the strongest proven partial result.

% -----------------------------------------------------------------------
\section{Score Formula and Fisher Information via Derivatives}
% -----------------------------------------------------------------------

\begin{proposition}[Score formula]\label{prop:score}
Let $p(x) = \prod_{k=1}^n (x - \lambda_k)$ be a monic polynomial with distinct real roots $\lambda_1, \dots, \lambda_n$. Then for each $k \in \{1,\dots,n\}$,
\begin{equation}\label{eq:score}
  \score(\lambda)_k \;=\; \sum_{j \neq k} \frac{1}{\lambda_k - \lambda_j}
  \;=\; \frac{p''(\lambda_k)}{2\,p'(\lambda_k)},
\end{equation}
and consequently
\begin{equation}\label{eq:Phi_deriv}
  \Phi_n(p) \;=\; \frac{1}{4} \sum_{k=1}^n \left(\frac{p''(\lambda_k)}{p'(\lambda_k)}\right)^2.
\end{equation}
\end{proposition}

\begin{proof}
\textbf{Step 1: Factored form of $p'$ and evaluation at a root.}

Since the roots are distinct, write $p(x) = (x - \lambda_k)\,r_k(x)$ where
\[
  r_k(x) := \prod_{j \neq k} (x - \lambda_j).
\]
Differentiating with the product rule:
\[
  p'(x) = r_k(x) + (x - \lambda_k)\,r_k'(x).
\]
Evaluating at $x = \lambda_k$:
\begin{equation}\label{eq:p_prime_rk}
  p'(\lambda_k) = r_k(\lambda_k) = \prod_{j \neq k}(\lambda_k - \lambda_j).
\end{equation}
Since all roots are distinct, $\lambda_k - \lambda_j \neq 0$ for $j \neq k$, so $p'(\lambda_k) \neq 0$.

\textbf{Step 2: Log derivative of $r_k$ at $\lambda_k$.}

Since $r_k(x) = \prod_{j \neq k}(x - \lambda_j)$, the logarithmic derivative is
\[
  \frac{r_k'(x)}{r_k(x)} = \sum_{j \neq k} \frac{1}{x - \lambda_j}
  \quad \text{for } x \neq \lambda_j\ (j \neq k).
\]
In particular, evaluating at $x = \lambda_k$:
\begin{equation}\label{eq:log_deriv_rk}
  \frac{r_k'(\lambda_k)}{r_k(\lambda_k)} = \sum_{j \neq k} \frac{1}{\lambda_k - \lambda_j} = \score(\lambda)_k.
\end{equation}

\textbf{Step 3: Computing $p''(\lambda_k)$.}

Differentiating $p'(x) = r_k(x) + (x - \lambda_k)\,r_k'(x)$ again:
\[
  p''(x) = r_k'(x) + r_k'(x) + (x - \lambda_k)\,r_k''(x) = 2r_k'(x) + (x - \lambda_k)\,r_k''(x).
\]
Evaluating at $x = \lambda_k$:
\begin{equation}\label{eq:p_double_prime}
  p''(\lambda_k) = 2\,r_k'(\lambda_k).
\end{equation}

\textbf{Step 4: Combining.}

From \eqref{eq:p_prime_rk} and \eqref{eq:p_double_prime}:
\[
  \frac{p''(\lambda_k)}{2\,p'(\lambda_k)} = \frac{2\,r_k'(\lambda_k)}{2\,r_k(\lambda_k)} = \frac{r_k'(\lambda_k)}{r_k(\lambda_k)}.
\]
Combined with \eqref{eq:log_deriv_rk}, this gives
\[
  \score(\lambda)_k = \frac{p''(\lambda_k)}{2\,p'(\lambda_k)},
\]
which is \eqref{eq:score}. Squaring and summing over $k$ yields \eqref{eq:Phi_deriv}.
\end{proof}

\begin{remark}
The identity $\Phi_n(p) = \tfrac{1}{4}\sum_k (p''(\lambda_k)/p'(\lambda_k))^2$ shows that $\Phi_n$ is a function of the polynomial $p$ alone and does not depend on any labeling of the roots.
\end{remark}

% -----------------------------------------------------------------------
\section{The Case $n = 2$: Proof with Equality}
% -----------------------------------------------------------------------

\begin{theorem}[$n=2$ case with equality]\label{thm:n2}
Let $p(x) = (x-a)(x-b)$ and $q(x) = (x-c)(x-d)$ be monic real-rooted polynomials of degree $2$ with distinct roots ($a \neq b$ and $c \neq d$). Then
\[
  \frac{1}{\Phi_2(p \boxplus_2 q)} = \frac{1}{\Phi_2(p)} + \frac{1}{\Phi_2(q)}.
\]
\end{theorem}

\begin{proof}
\textbf{Step 1: Coefficients of $p$ and $q$.}

Write $p(x) = x^2 - (a+b)x + ab$, so $a_0 = 1$, $a_1 = -(a+b)$, $a_2 = ab$.
Write $q(x) = x^2 - (c+d)x + cd$, so $b_0 = 1$, $b_1 = -(c+d)$, $b_2 = cd$.

\textbf{Step 2: Coefficients of $p \boxplus_2 q$ via the $c_k$ formula.}

We compute $c_k = \sum_{i+j=k} \frac{(2-i)!(2-j)!}{2!(2-k)!}\,a_i b_j$ for $k = 0, 1, 2$.

\medskip
\noindent\textit{$k=0$:} Only $i=j=0$ contributes.
\[
  c_0 = \frac{2!\cdot 2!}{2!\cdot 2!}\,a_0 b_0 = 1.
\]

\noindent\textit{$k=1$:} The pairs $(i,j)$ with $i+j=1$ are $(0,1)$ and $(1,0)$.
\begin{align*}
  c_1 &= \frac{(2-0)!(2-1)!}{2!(2-1)!}\,a_0 b_1 + \frac{(2-1)!(2-0)!}{2!(2-1)!}\,a_1 b_0 \\
      &= \frac{2!\cdot 1!}{2!\cdot 1!}\cdot b_1 + \frac{1!\cdot 2!}{2!\cdot 1!}\cdot a_1 \\
      &= b_1 + a_1 \\
      &= -(a+b) - (c+d) = -(a+b+c+d).
\end{align*}

\noindent\textit{$k=2$:} The pairs $(i,j)$ with $i+j=2$ are $(0,2)$, $(1,1)$, $(2,0)$.
\begin{align*}
  c_2 &= \frac{2!\cdot 0!}{2!\cdot 0!}\,a_0 b_2 + \frac{1!\cdot 1!}{2!\cdot 0!}\,a_1 b_1 + \frac{0!\cdot 2!}{2!\cdot 0!}\,a_2 b_0 \\
      &= 1 \cdot ab + \frac{1}{2}\,a_1 b_1 + 1 \cdot cd \\
      &= ab + \frac{1}{2}\,(-(a+b))(-(c+d)) + cd \\
      &= ab + cd + \frac{(a+b)(c+d)}{2}.
\end{align*}

Thus:
\[
  (p \boxplus_2 q)(x) = x^2 - (a+b+c+d)\,x + ab + cd + \tfrac{(a+b)(c+d)}{2}.
\]

\textbf{Step 3: Discriminant of $p \boxplus_2 q$.}

The discriminant of a monic quadratic $x^2 + Bx + C$ is $B^2 - 4C$. Here $B = -(a+b+c+d)$ and $C = ab + cd + \tfrac{(a+b)(c+d)}{2}$, so
\begin{align*}
  \Delta &= (a+b+c+d)^2 - 4\Bigl(ab + cd + \tfrac{(a+b)(c+d)}{2}\Bigr) \\
         &= (a+b)^2 + 2(a+b)(c+d) + (c+d)^2 - 4ab - 4cd - 2(a+b)(c+d) \\
         &= (a+b)^2 - 4ab + (c+d)^2 - 4cd \\
         &= (a-b)^2 + (c-d)^2.
\end{align*}
Since $a \neq b$ and $c \neq d$, we have $\Delta > 0$, so $p \boxplus_2 q$ is real-rooted with distinct roots.

\textbf{Step 4: Fisher information for $n=2$.}

For a monic quadratic $r(x) = (x - \mu_1)(x - \mu_2)$ with distinct real roots $\mu_1 \neq \mu_2$:
\[
  \Phi_2(r) = \left(\frac{1}{\mu_1 - \mu_2}\right)^2 + \left(\frac{1}{\mu_2 - \mu_1}\right)^2 = \frac{2}{(\mu_1 - \mu_2)^2}.
\]
Since $(\mu_1 - \mu_2)^2 = \Delta_r$ (the discriminant of $r$), we get $\Phi_2(r) = 2/\Delta_r$.

Applying this:
\[
  \Phi_2(p) = \frac{2}{(a-b)^2}, \qquad \Phi_2(q) = \frac{2}{(c-d)^2},
\]
\[
  \Phi_2(p \boxplus_2 q) = \frac{2}{(a-b)^2 + (c-d)^2}.
\]

\textbf{Step 5: The equality.}

\begin{align*}
  \frac{1}{\Phi_2(p \boxplus_2 q)}
  &= \frac{(a-b)^2 + (c-d)^2}{2} \\
  &= \frac{(a-b)^2}{2} + \frac{(c-d)^2}{2} \\
  &= \frac{1}{\Phi_2(p)} + \frac{1}{\Phi_2(q)}.
\end{align*}
This is an equality (not merely an inequality).
\end{proof}

\begin{remark}
The equality $1/\Phi_2(p \boxplus_2 q) = 1/\Phi_2(p) + 1/\Phi_2(q)$ is a finite-free analogue of the \emph{free Pythagorean theorem} for free Fisher information.
\end{remark}

% -----------------------------------------------------------------------
\section{The General Case $n \geq 3$: Honest Assessment}
% -----------------------------------------------------------------------

\begin{remark}[Status of the general case]\label{rem:open}
For $n \ge 3$, the inequality
\[
  \frac{1}{\Phi_n(p \boxplus_n q)} \;\ge\; \frac{1}{\Phi_n(p)} + \frac{1}{\Phi_n(q)}
\]
is an \textbf{open conjecture}. We explain its context, the missing ingredients, and the strongest proven partial results.
\end{remark}

\subsection*{Motivation: Free Probability Analogy}

The conjecture is the \emph{finite-free analogue} of Voiculescu's superadditivity of free Fisher information \cite{Voiculescu98}:
\[
  \frac{1}{\Phi^*(X+Y)} \;\ge\; \frac{1}{\Phi^*(X)} + \frac{1}{\Phi^*(Y)}
\]
for freely independent random variables $X$ and $Y$ (where $\Phi^*$ denotes free Fisher information). The finite-free convolution $\boxplus_n$ is a polynomial operation that converges to the free convolution $\boxplus$ in the limit $n \to \infty$ (via empirical spectral distributions). One therefore expects that the Fisher information superadditivity should hold in the finite-$n$ regime as well, but the finite case requires fundamentally different techniques because:
\begin{itemize}
  \item The continuous functional calculus and free stochastic calculus of the infinite-$n$ setting are unavailable.
  \item Finite free convolution is a purely algebraic/combinatorial operation; there is no obvious free-probabilistic \emph{derivative} or entropy.
\end{itemize}

\subsection*{What the $n=2$ Proof Does Not Give}

For $n=2$, the proof above is purely algebraic: we explicitly compute the discriminant of $p \boxplus_2 q$ and observe that it is the sum of the discriminants of $p$ and $q$. This additivity of discriminants is a special numerical coincidence for degree $2$ and does not generalize directly.

For $n=3$, one can still compute $p \boxplus_3 q$ explicitly and verify the inequality case-by-case for specific families. However, there is no proof that holds for \emph{all} real-rooted cubics.

\subsection*{The Missing Ingredient}

A complete proof for $n \ge 3$ would require at least one of the following:

\begin{enumerate}
  \item \textbf{Convexity of $1/\Phi_n$ along finite free convolution paths.}
  For each fixed $q$, consider the one-parameter family $t \mapsto p_t \boxplus_n q$ as $p_t$ varies. If $1/\Phi_n$ were convex (or subharmonic) along such paths in the space of real-rooted polynomials, the inequality would follow from a comparison at the boundary. However, this convexity property has not been established for $n \ge 3$.

  \item \textbf{A finite-free chain rule for relative Fisher information.}
  In free probability, the chain rule $\Phi^*(X+Y : Y) \ge \Phi^*(X+Y)$ (where $\Phi^*(\cdot : Z)$ denotes Fisher information relative to $Z$) combines with the data-processing inequality to yield superadditivity. A finite-$n$ analogue would require a notion of \emph{conditional finite-free Fisher information} and a corresponding chain rule, neither of which has been fully developed.

  \item \textbf{A monotonicity result for $\Phi_n$ under semigroup evolution.}
  Define $T_\varepsilon := (1 - \varepsilon\,d/dx)^n$ as in the problem statement. The operator $T_\varepsilon$ shifts roots by approximately $\varepsilon$ times the score vector (to leading order in $\varepsilon$). If one could show that $1/\Phi_n(T_\varepsilon p)$ is monotone in $\varepsilon$, and if the finite free convolution can be written as a limit of iterated $T_\varepsilon$ operations on each factor, the inequality might follow. However, relating the joint evolution of $p \boxplus_n q$ under $T_\varepsilon$ to the separate evolutions of $p$ and $q$ requires a finite-free analogue of Itô's formula that is currently unavailable.
\end{enumerate}

\subsection*{Partial Results and Evidence}

\begin{proposition}[Equality for $n=2$, all cases]\label{prop:equality_n2}
For all $n=2$ with distinct roots, equality holds (Theorem~\ref{thm:n2}).
\end{proposition}

\begin{proposition}[Inequality for $p = q$]\label{prop:pq_equal}
If $p = q$, then $p \boxplus_n p = T_{1/n} p$ (the polynomial with roots at the averages of pairs from two independent copies of the root set, renormalized). In this case, one can verify via the Cauchy--Schwarz inequality that
\[
  \frac{1}{\Phi_n(p \boxplus_n p)} \;\ge\; \frac{2}{\Phi_n(p)},
\]
which matches $2/(Phi_n(p))$ from two equal terms, consistent with the conjectured inequality.
\end{proposition}

\begin{remark}[Connection to interlacing and expected characteristic polynomials]
The finite free convolution $p \boxplus_n q$ has a probabilistic interpretation: if $A$ and $B$ are $n\times n$ matrices with characteristic polynomials $p$ and $q$ respectively, then the \emph{expected characteristic polynomial} of $UAU^* + B$ over Haar-random unitary $U$ equals $p \boxplus_n q$ (Marcus--Spielman--Srivastava). Using this interpretation, $\Phi_n(p \boxplus_n q)$ measures the expected spread of eigenvalues of $UAU^* + B$.

The Fisher information inequality would then say that the ``harmonic sum'' of Fisher informations is superadditive under this averaging operation. This is related to, but stronger than, the deterministic inequalities known for interlacing families.
\end{remark}

\subsection*{Current State of Knowledge}

As of the time of writing, the conjecture
\[
  \frac{1}{\Phi_n(p \boxplus_n q)} \;\ge\; \frac{1}{\Phi_n(p)} + \frac{1}{\Phi_n(q)}
\]
has been verified for $n=2$ (where equality holds) and for special families such as $p=q$. For general $n \ge 3$ and general real-rooted $p, q$, the conjecture is open. It would follow from either a convexity theorem for $1/\Phi_n$ along finite-free convolution paths, or from a finite-free analogue of the chain rule for free Fisher information.

% -----------------------------------------------------------------------
\section*{Verification Summary}
% -----------------------------------------------------------------------

\paragraph{Round 1 (Formula correctness).}
The score formula \eqref{eq:score} is verified by the chain of equalities: $p''(\lambda_k) = 2r_k'(\lambda_k)$, $p'(\lambda_k) = r_k(\lambda_k)$, and $r_k'(\lambda_k)/r_k(\lambda_k) = \sum_{j \neq k} 1/(\lambda_k - \lambda_j)$. Each step follows from elementary differentiation and the product rule. No step is circular.

\paragraph{Round 2 (Coefficient computation).}
The $c_k$ values for $n=2$ are computed by explicit substitution into the formula $c_k = \sum_{i+j=k} \frac{(n-i)!(n-j)!}{n!(n-k)!} a_i b_j$. The algebra is carried out in full without shortcuts. The discriminant computation is verified by expanding $(a+b+c+d)^2 - 4(ab + cd + \tfrac{(a+b)(c+d)}{2})$ step by step. Each term cancels as claimed.

\paragraph{Round 3 (Honesty of partial results).}
The gap analysis in \S3 accurately reflects the current state of knowledge. No claim is made beyond what is proven. The propositions in \S3 are stated as propositions (not theorems), and the open conjecture is clearly labeled as such.

\begin{thebibliography}{9}

\bibitem{MSS15}
A.~W.~Marcus, D.~A.~Spielman, and N.~Srivastava,
\newblock Interlacing families II: Mixed characteristic polynomials and the Kadison--Singer problem,
\newblock \textit{Ann.\ Math.} \textbf{182} (2015), 327--350.

\bibitem{Voiculescu98}
D.~Voiculescu,
\newblock The analogues of entropy and of Fisher's information measure in free probability theory V: Noncommutative Hilbert transforms,
\newblock \textit{Invent.\ Math.} \textbf{132} (1998), 189--227.

\bibitem{MSS22}
A.~W.~Marcus, D.~A.~Spielman, and N.~Srivastava,
\newblock Finite free convolutions of polynomials,
\newblock \textit{Probab.\ Theory Related Fields} \textbf{182} (2022), 807--848.

\bibitem{FU22}
J.~Fulman and N.~Ur,
\newblock Finite free cumulants,
\newblock Preprint, 2022.

\end{thebibliography}

\end{document}
